\contourlength{1.4pt}

\newcommand{\calI}{\mathscr{I}} %\mathcal
\tikzset{>=latex} % for LaTeX arrow head
\colorlet{myred}{red!80!black}
\colorlet{myblue}{blue!80!black}
\colorlet{mygreen}{green!80!black}
\colorlet{mydarkred}{red!50!black}
\colorlet{mydarkblue}{blue!50!black}
\colorlet{mylightblue}{mydarkblue!6}
\colorlet{mypurple}{blue!40!red!80!black}
\colorlet{mydarkpurple}{blue!40!red!50!black}
\colorlet{mylightpurple}{mydarkpurple!80!red!6}
\colorlet{myorange}{orange!40!yellow!95!black}
\tikzstyle{cone}=[mydarkblue,line width=0.2,top color=blue!60!black!30,
                  bottom color=blue!60!black!50!red!30,shading angle=60,fill opacity=0.9]
\tikzstyle{cone back}=[mydarkblue,line width=0.1,dash pattern=on 1pt off 1pt]
\tikzstyle{world line}=[myblue!60,line width=0.4]
\tikzstyle{world line t}=[mypurple!60,line width=0.4]
\tikzstyle{particle}=[mygreen,line width=0.5]
\tikzstyle{photon}=[-{Latex[length=4,width=3]},myorange,line width=0.4,decorate,
                    decoration={snake,amplitude=0.9,segment length=4,post length=3.8}]
\tikzstyle{singularity}=[myred,line width=0.6,decorate,
                         decoration={zigzag,amplitude=2,segment length=6.17}]
\tikzset{declare function={%
  penrose(\x,\c)  = {\fpeval{2/pi*atan( (sqrt((1+tan(\x)^2)^2+4*\c*\c*tan(\x)^2)-1-tan(\x)^2) /(2*\c*tan(\x)^2) )}};%
  penroseu(\x,\t) = {\fpeval{atan(\x+\t)/pi+atan(\x-\t)/pi}};%
  penrosev(\x,\t) = {\fpeval{atan(\x+\t)/pi-atan(\x-\t)/pi}};%
  kruskal(\x,\c)  = {\fpeval{asin( \c*sin(2*\x) )*2/pi}};% Penrose coordinates for Kruskal
}}
\def\tick#1#2{\draw[thick] (#1) ++ (#2:0.04) --++ (#2-180:0.08)}
\def\Nsamples{20} % number samples in plot

\def\R{0.08} % size lightcone
\def\e{0.08} % vertical scale
\def\ang{45} % angle light cone
\def\angb{acos(sqrt(\e)*sin(\ang))} % angle ellipse center to point of tangency
\def\a{\R*sin(\ang)*sqrt(1-\e*sin(\ang)^2)/(1-\e*sin(\ang)^2)} % vertical radius
\def\b{\R*sqrt(\e)*sin(\ang)*cos(\ang)/(1-\e*sin(\ang)^2)} % horizontal radius
\def\coneback#1{ % light cone part to be drawn behind world lines
  \draw[cone back] % dashed line back
    (#1)++(-45:\R) arc({90-\angb}:{90+\angb}:{\a} and {\b});
  \draw[cone,shading angle=-60] % top edge & inside
    (#1)++(0,{\R*cos(\ang)/(1-\e*sin(\ang)^2)}) ellipse({\a} and {\b});
}
\def\conefront#1{ % light cone part to be drawn over world lines
  \draw[cone] % light cone outside
    (#1) --++ (45:\R) arc({\angb-90}:{-90-\angb}:{\a} and {\b})
     --++ (-45:2*\R) arc({90-\angb}:{-270+\angb}:{\a} and {\b}) -- cycle;
}

% PENROSE DIAGRAM of Minkowski space - 45 rotation




% PENROSE DIAGRAM of Minkowski space - equidistant world lines


\begin{tikzpicture}[scale=3.2]
  \message{Penrose diagram (equidistant world lines)^^J}
  
  \def\Nlines{9} % number of world lines (at constant r/t)
  \def\d{0.5}
  \coordinate (O) at ( 0, 0); % center: origin (r,t) = (0,0)
  \coordinate (S) at ( 0,-1); % south: t=-infty, i-
  \coordinate (N) at ( 0, 1); % north: t=+infty, i+
  \coordinate (W) at (-1, 0); % east:  r=-infty, i0
  \coordinate (E) at ( 1, 0); % west:  r=+infty, i0
  \coordinate (X) at ({penroseu(\d,\d)},{penrosev(\d,\d)});
  \coordinate (X0) at ({penroseu(\d,-2*\d)},{penrosev(\d,-2*\d)});
  
  % AXES
  \fill[mylightblue] (N) -- (E) -- (S) -- (W) -- cycle;
  \draw[->,thick] (-1.1,0) -- (1.15,0) node[right=-1] {$u$};
  \draw[->,thick] (0,-1.1) -- (0,1.15) node[left=-1] {$v$};
  
  % INFINITY LABELS
  \node[right=1,above=1,mydarkblue] at (-1,0.04) {$i^0$};
  \node[right=1,above=1,mydarkblue] at (1,0.04) {$i^0$};
  \node[below=0,right=1,mydarkpurple] at (0.04,-1) {$i^-$};
  \node[above=1,right=1,mydarkpurple] at (0.04,1) {$i^+$};
  \node[mydarkblue,below left=-1] at (-0.5,-0.5) {$\calI^-$};
  \node[mydarkblue,above left=-1] at (-0.5,0.5) {$\calI^+$};
  \node[mydarkblue,above right=-1] at (0.5,0.5) {$\calI^+$};
  \node[mydarkblue,below right=-1] at (0.5,-0.5) {$\calI^-$};
  
  % LIGHT CONE BACK
  \coneback{X};
  \coneback{X0};
  
  % WORLD LINES
  \draw[world line] (N) -- (S);
  \draw[world line t] (W) -- (E);
  \message{Making world lines...^^J}
  \foreach \i [evaluate={\c=\i*\d;}] in {1,...,\Nlines}{
    \message{  Running i/N=\i/\Nlines, c=\c...^^J}
    \draw[world line t,samples=\Nsamples,smooth,variable=\x,domain=-1:1] % constant t
      plot(\x,{-penrose(\x*pi/2,\c)})
      plot(\x,{ penrose(\x*pi/2,\c)});
    \draw[world line,samples=\Nsamples,smooth,variable=\y,domain=-1:1] % constant r
      plot({-penrose(\y*pi/2,\c)},\y)
      plot({ penrose(\y*pi/2,\c)},\y);
  }
  \draw[thick,mydarkblue] (N) -- (E) -- (S) -- (W) -- cycle;
  
  % PARTICLE
  \draw[particle,decoration={markings,mark=at position 0.16 with {\arrow{latex}},
                                      mark=at position 0.53 with {\arrow{latex}},
                                      mark=at position 0.81 with {\arrow{latex}}},postaction={decorate}]
    (S) to[out=90,in=-80] (X0) to[out=100,in=-95] (X) to[out=85,in=-90] (N);
  
  % LIGHT CONE FRONT
  \conefront{X};
  \conefront{X0};
  
  % PHOTON
  \draw[->,photon] (O) -- (-0.5,0.5); %node[above=2,right=1] {photon};
  
  % TICKS
  \tick{W}{90} node[left=4,below=-1] {$-\pi/2$};
  \tick{E}{90} node[right=4,below=-1] {$+\pi/2$};
  \tick{S}{ 0} node[left=-1] {$-\pi/2$};
  \tick{N}{ 0} node[left=-1] {$+\pi/2$};
  
\end{tikzpicture}
